\chapter{Tensor semistability}

\begin{lemma}[Degree criterion and normalization]
\label{lem:semistable-normalization}
\uses{lem:degree-calculus}
\source{bost-specht-manuscript:page-0012}
\lean{BostConjecture.Arithmetic.HermitianMetric.isSemistable_iff_degree_nonpos,BostConjecture.Arithmetic.HermitianMetric.isSemistable_rescale_iff,BostConjecture.Arithmetic.HermitianMetric.degree_rescale_self,BostConjecture.Arithmetic.HermitianMetric.slope_submodule_rescale}
\leanok
For a Hermitian bundle $E$ of degree zero, semistability is equivalent
to $\degn(B)\leq0$ for every integral submodule $B\subset E$,
including zero. If $E$ has positive rank, multiplying every norm by
$e^t$ preserves semistability and subtracts $t$ from every
positive-rank slope. Taking $t=\slope(E)$ normalizes its degree.
\end{lemma}
\begin{proof}
The slope inequality for a positive-rank submodule is its degree
inequality divided by its positive rank; a rank-zero submodule has
degree zero. Uniform scaling changes the determinant norm on a
rank-$b$ submodule by $e^{bt}$ at every embedding. The normalized
degree therefore decreases by $bt$, so its slope decreases by $t$.
The ambient slope changes by the same amount, preserving every
comparison. The specified parameter makes the ambient slope and
hence degree zero.
\end{proof}

\begin{lemma}[Normalized extensions and degree-zero pieces]
\label{lem:normalized-semistable-extension}
\uses{lem:degree-calculus,lem:semistable-normalization}
\source{bost-specht-manuscript:page-0012,bost-specht-manuscript:page-0013}
\lean{BostConjecture.Arithmetic.HermitianMetric.isSemistable_of_submodule_quotient,BostConjecture.Arithmetic.HermitianMetric.IsSemistable.submodule_of_degree_zero,BostConjecture.Arithmetic.HermitianMetric.IsSemistable.quotient_of_degree_zero}
\leanok
Let $A$ be a saturated submodule of a Hermitian bundle $E$ and equip
$A$ and $Q=E/A$ with induced and quotient metrics.
If $A,Q$ are semistable of degree zero, then so is $E$.
Conversely, if $E$ is semistable of degree zero and $\degn(A)=0$,
then $A,Q$ are both semistable of degree zero.
\end{lemma}
\begin{proof}
For the first assertion, exact degree gives $\degn(E)=0$.
For any $S\subset E$, put $I=(S+A)/A\subset Q$.
The sum/intersection inequality and exact quotient degree give
\[
 \degn(S)\leq\degn(S\cap A)+\degn(S+A)-\degn(A)
 =\degn(S\cap A)+\degn(I)\leq0.
\]
The two final bounds follow from semistability of $A,Q$, also when
one of the modules has rank zero. The degree criterion proves
semistability of $E$.
For the converse, every submodule of $A$ has the same induced metric
when viewed in $E$, so it has nonpositive degree. Exact degree gives
$\degn(Q)=0$. For any submodule $I\subset Q$, its inverse image
$T\subset E$ satisfies $\degn(I)=\degn(T)-\degn(A)\leq0$.
The degree criterion applied to $A$ and $Q$ concludes the proof.
\end{proof}

\begin{lemma}[Extensions and line twists]
\label{lem:extensions-and-twists}\uses{lem:degree-calculus,lem:normalized-semistable-extension}\notready
An extension of normalized semistable bundles, with induced and
quotient metrics, is normalized semistable. Tensoring by a Hermitian
line bundle preserves semistability and shifts every slope by the
line's degree.
\end{lemma}
\begin{proof}
For $0\to A\to E\to Q\to0$ and a nonzero submodule $S\subset E$,
let $B=S\cap A$ and let $I$ be its actual image in $Q$. The
isomorphism $S/B\to I$ is contracting, because distance to $A$ is
no greater than distance to $B$. Degree additivity and contraction
therefore give
$\degn(S)\leq\degn(B)+\degn(I)\leq0$.
Also $\degn(E)=\degn(A)+\degn(Q)=0$, proving the first assertion.

An invertible module $L$ gives inverse correspondences
$B\mapsto L\otimes B$ and $C\mapsto L^\vee\otimes C$ between
submodules of $E$ and of $L\otimes E$. Local trivializations show
that ranks and saturation are preserved. At each archimedean place,
the line tensor factor multiplies every vector norm by the norm of a
line generator; hence these correspondences preserve the induced
metric up to that same line twist. The degree formula shifts every
submodule slope, as well as the ambient slope, by $\degn(L)$, so
their comparisons are unchanged.
\end{proof}

\begin{proposition}[Tensor semistability]
\label{thm:tensor-semistability}
\uses{lem:degree-calculus,lem:extensions-and-twists,lem:arithmetic-compactness,lem:degree-maximum-continuity,lem:balanced-plucker,lem:binary-plucker-nonzero-sequence,thm:geometric-slope}
\source{bost-specht-manuscript:page-0012,bost-specht-manuscript:page-0013}
\notready
For every number field $K$ and all positive-rank semistable Hermitian
$\Ocal_K$-bundles $E,F$, their Hilbert tensor product is semistable.
\end{proposition}
\begin{proof}
Uniformly scale the norms on each factor to make its slope zero.
This does not change semistability, and changes every slope of the
tensor product by the sum of the two shifts. Suppose a normalized
counterexample exists and choose one with $m+n$ minimal, where
$m=\rk E$ and $n=\rk F$. The line-twist lemma implies $m,n\geq2$.

The function $\mathcal D(E\otimes F)$ is bounded above and attains
its supremum on normalized semistable pairs of these ranks. Indeed,
otherwise a sequence of values tending to infinity, or to an unattained
finite supremum, would by arithmetic compactness have a subsequence
with fixed modules and convergent metrics; tensor metrics converge
and continuity of $\mathcal D$ gives the corresponding finite limiting
value, a contradiction in the first case and attainment in the second.
Choose a maximizing pair and a saturated submodule $S\subset E\otimes F$
of maximal degree. Since a counterexample exists, $\degn(S)>0$.

If $E$ had a nonzero proper saturated submodule $A$ of degree zero,
then $A$ is normalized semistable by restriction of the inequalities
for $E$. Its quotient $E/A$ is normalized semistable too: the inverse
image of a submodule of positive degree would, by degree additivity,
have positive degree in $E$. By minimality of $m+n$, both
$A\otimes F$ and $(E/A)\otimes F$ are normalized semistable. Tensoring
the exact sequence by the projective module $F$ preserves exactness.
Its induced and quotient metrics agree with the Hilbert tensor
metrics, by an orthogonal decomposition of $E$ along $A$.
The extension lemma makes $E\otimes F$ semistable, a contradiction.
Thus $E$ is strictly stable. The same reasoning proves strict stability
of $F$.

Fix an archimedean place $v$. By openness of strict stability, small
determinant-preserving variations of the metrics at $v$ remain
normalized semistable; at complex places vary its conjugate as well.
The degree of the fixed module $S$ has a local maximum at the chosen
pair, since it is bounded above by $\mathcal D$ for each nearby pair,
and that function is bounded above by its chosen global maximum.
Let $P$ be the orthogonal projection onto $S_v$ and $s=\rk S$.
For a traceless Hermitian endomorphism $H$ of $E_v$, change its norm
to $x\mapsto\|e^{tH}x\|$. Differentiating the determinant norm of
$S_v$, using an orthonormal basis of $S_v$, gives
\[
 \left.\frac{d}{dt}\degn(S)\right|_{t=0}
 =-\frac{w_v}{D}\Tr(P(H\otimes\Id))=0.
\]
Consequently $\Tr_F P$ is orthogonal under the trace pairing to all
traceless Hermitian matrices. Taking
$H=\Tr_F P-(\Tr P/m)\Id$ shows that the trace of its square is
zero; since it is Hermitian, it is zero. As $\Tr P=s$, this gives
$\Tr_F P=(s/m)\Id$. Varying $F$ similarly gives
$\Tr_E P=(s/n)\Id$. At a real place the same argument uses real
symmetric matrices, and its resulting matrix identities persist after
complexification.

The balanced-projection lemma supplies a positive-degree homogeneous
Pl\"ucker invariant nonzero on the exterior vector of $S$ at this complex
embedding. Binary invariant detection gives nonzero canonical contractions
over $\Ocal_K$ in every positive multiple of a common degree. The contraction
and Specht-asymptotic argument of the geometric slope inequality therefore
gives $\slope(S)\leq\slope(E)+\slope(F)=0$, contradicting
$\degn(S)>0$. No counterexample exists.
\end{proof}

